\section{The comparison target and the meaning of simplicity}

The reference is the exact repository snapshot displayed on the title
page, inspected through GitHub. Its main native kernel $\mathcal N$
has equations, a one-variable closure $Q$, and a two-variable binder $D$.
It has eight primitive rule constructors: hypothesis, two computation
equations, equational transport, $Q$ elimination, $D$ elimination,
$D$ introduction by cases, and $Q$ introduction by equation progress.
The repository proves its equivalence with full first-order iPRA in Lean
\cite{baseline,native,bridge}.

The reference also contains a seven-rule presentation $W$.
Its single introduction uses progress exactly when the matrix is an
equation and the witness is the literal zero; it uses cases otherwise.
The theorem \texttt{wprf\_iff\_nprf} proves that this presentation has the
same sequents as the native kernel. Consequently, a claim of novelty
based on the number seven alone would be incorrect. A numerical
substitution can change a witness into literal zero, so a direct induction
on $W$ derivations does not preserve the primitive introduction choice
\cite{witness}.

The present construction changes the meaning and elimination of atomic
closures. This makes the introduction choice depend only on the formula
root, which numerical substitution preserves. Unlike the guarded
presentations \cite{guarded}, every formula generated by the displayed
grammar is public. No restriction to a prescribed zero guard, protected
equation, or mirror shape is imposed.

\begin{center}
\small
\begin{tabular}{@{}p{.18\textwidth}p{.18\textwidth}p{.29\textwidth}p{.24\textwidth}@{}}
\toprule
Presentation & Introductions & Primitive mode selection & Public formulas \\
\midrule
Adopted $\mathcal N$ & Cases and equation progress & Two rules & All Eq/$Q$/$D$ \\
Known $W$ & One & Matrix root and literal witness & All Eq/$Q$/$D$ \\
Guarded core & One & Matrix root & Restricted atomic matrices \\
New $\mathcal R$ & One & Matrix root only & All Eq/$Q$/$D$ \\
\bottomrule
\end{tabular}
\end{center}

The scalar terms, arithmetic construction, and use of progressive
conditions are inherited ideas. In particular, an earlier progressive
presentation uses progress at every matrix root, counter-aware recursion,
separate step and iteration eliminations, and an explicit Boolean-cases
rule \cite{progressive}. Here the iterator remains scalar, compound
closures retain their conditional reading, there is no primitive
Boolean-cases rule, and one orbit specialization serves both kinds of
closure. The all-formula normalization theorem is the substantive new
result.

\subsection{What the size claim is}

The relevant implementation criterion is the size of a complete
standalone proof checker, with gzip level 9 reported separately for
information. We measure significant Python tokens, all AST nodes, and
UTF-8 bytes after removing genuine docstrings and applying
\texttt{ast.unparse}, followed by one newline. No producer or theorem
database is imported by a checker.

Two comparisons are needed. The primary comparison uses explicit tagged
syntax, exact integer/list checks, de Bruijn scoping, and complete
Eq/$Q$/$D$ statement validation in three matched standalone programs.
The secondary comparison concerns the repository's much smaller arity
codec under its canonical-public-goal convention. That convention does
not include a complete public-grammar validator. Its numbers cannot be
interchanged with those of the typed interface.

The inspected incoming archives are part of the comparison, not omitted
competitors. Their original ZIP bytes were authenticated against the six
Git blob identifiers at the pinned commit. They include
\texttt{deterministic\_auxiliary}, \texttt{guarded\_source},
\texttt{guarded\_orbit}, \texttt{mirror\_guard},
\texttt{protected\_core}, and \texttt{reversal\_core}
\cite{incoming}. Exact source records and the comparison measurements
are included in the archive.

\textbf{Scope of the outcome.}
The new theory gives a cleaner introduction discipline and preserves full
bi-interpretability. It does not improve every existing checker-size
coordinate. The separate compact implementation gain is small and does
not establish a new mathematical lower bound. No global minimality
theorem is claimed.
